\section{El Launch File}

Un launch file es un script Python cuya única responsabilidad es describir qué nodos lanzar, con qué parámetros y en qué orden. \texttt{ros2 launch} lo importa como módulo y busca la función \texttt{generate\_launch\_description()}, que debe retornar un objeto \texttt{LaunchDescription}: una lista ordenada de acciones que el sistema procesa en secuencia.

Toda la configuración de parámetros pasa por el mismo par de piezas: \texttt{DeclareLaunchArgument} registra el argumento y establece su valor por defecto; \texttt{LaunchConfiguration} lo recupera después. El valor \textit{siempre} viaja como string —ROS~2 nunca convierte tipos automáticamente.

Existen dos variantes del patrón según si se necesita o no ejecutar lógica Python sobre los valores.

% ------------------------------------------------------------------
\subsection{Caso 1: parámetros escalares simples}

Cuando los parámetros se pasan directamente al nodo sin transformación, \texttt{LaunchConfiguration} puede usarse como sustitución en línea. ROS~2 resuelve su valor en el momento de lanzar el proceso.

\begin{figure}[H]
\centering
\begin{lstlisting}[style=python,
    caption={Launch file sin lógica Python. \texttt{LaunchConfiguration}
             actúa como sustitución diferida.},
    label={lst:launch_simple}]
from launch import LaunchDescription
from launch.actions import DeclareLaunchArgument
from launch.substitutions import LaunchConfiguration
from launch_ros.actions import Node

def generate_launch_description():
    return LaunchDescription([
        DeclareLaunchArgument('vmax', default_value='0.5'),
        DeclareLaunchArgument('wmax', default_value='1.0'),

        Node(
            package='control_nodes',
            executable='consenso_node',
            namespace='robot0',
            parameters=[{
                'vmax': LaunchConfiguration('vmax'),
                'wmax': LaunchConfiguration('wmax'),
            }],
        ),
    ])
\end{lstlisting}
\end{figure}

Este patrón es suficiente cuando no hay cálculo ni conversión de tipos: ROS~2 entrega el string al nodo y el nodo lo interpreta.

% ------------------------------------------------------------------
\subsection{Caso 2: lógica Python sobre los argumentos}

Cuando se necesita procesar el valor en Python —convertir tipos, iterar, calcular retardos— se añade \texttt{OpaqueFunction}. Al llegar a esa acción, el sistema llama a la función indicada y le pasa el \textit{context}: el estado global de la ejecución que contiene todos los valores ya resueltos. Dentro de la función, \texttt{.perform(context)} extrae el string del context para que el código Python pueda operarlo.

\begin{figure}[H]
\centering
\begin{lstlisting}[style=python,
    caption={Launch file con \texttt{OpaqueFunction}. Se usa cuando
             hay conversión de tipos o cálculos sobre los argumentos.},
    label={lst:launch_opaque}]
from launch import LaunchDescription
from launch.actions import DeclareLaunchArgument, OpaqueFunction
from launch.substitutions import LaunchConfiguration
from launch_ros.actions import Node

def launch_setup(context, *args, **kwargs):
    vmax = float(LaunchConfiguration('vmax').perform(context))
    # Los waypoints llegan como string: "1.0,0.0,4.0,0.0"
    raw  = LaunchConfiguration('waypoints').perform(context)
    wps  = [float(v) for v in raw.split(',')]

    return [
        Node(
            package='control_nodes',
            executable='navigate_waypoints_pva',
            namespace='robot0',
            parameters=[{
                'vmax':      vmax,
                'waypoints': wps,
            }],
        )
    ]

def generate_launch_description():
    return LaunchDescription([
        DeclareLaunchArgument('vmax',
            default_value='0.5'),
        DeclareLaunchArgument('waypoints',
            default_value='4.0,0.0,4.0,4.0,0.0,4.0'),
        OpaqueFunction(function=launch_setup),
    ])
\end{lstlisting}
\end{figure}

La diferencia práctica: en el Caso~1, \texttt{LaunchConfiguration} es una promesa que ROS~2 resuelve solo; en el Caso~2, \texttt{.perform(context)} fuerza la resolución de inmediato para que el código Python pueda operar sobre el valor real.

% ------------------------------------------------------------------
\subsection{Flujo de un valor de extremo a extremo}

El tipo del argumento cambia solo donde el usuario escribe una conversión explícita: \texttt{int()} o \texttt{float()} en el launch file, e implícitamente al deserializar YAML en el Parameter Server. ROS~2 nunca convierte tipos por cuenta propia en ninguna de las etapas.
